\section{Vision transformer: overview} 
\label{sec:vit} 

In this section, we briefly recall preliminaries associated with the vision transformer~\cite{dosovitskiy2020image,Vaswani2017AttentionIA}, and further discuss positional encoding and resolution. 

%\subsection{Visual transformers: overview}

\paragraph{Multi-head Self Attention layers (MSA).}
The attention mechanism is based on a trainable associative memory with (key, value) vector pairs. 
% 
A \emph{query} vector $q\in\mathbb{R}^d$ is matched against a set of $k$ \emph{key} vectors (packed together into a matrix $K\in\mathbb{R}^{k\times d}$) using inner products. 
These inner products are then scaled and normalized with a softmax function to obtain $k$ weights. 
The output of the attention is the weighted sum of a set of $k$ \emph{value} vectors (packed into $V\in\mathbb{R}^{k\times d}$). 
For a sequence of $N$ query vectors (packed into $Q\in\mathbb{R}^{N\times d}$), it produces an output matrix (of size $N\times d$): 
\begin{equation}
    \mathrm{Attention}(Q, K, V) = \mathrm{Softmax}(Q K^\top/\sqrt{d}) V,
\end{equation}
where the $\mathrm{Softmax}$ function is applied over each row of the input matrix and the $\sqrt{d}$ term provides appropriate normalization.\\
In~\cite{Vaswani2017AttentionIA}, a Self-attention layer is proposed.  Query, key and values matrices are themselves computed from a sequence of $N$ input vectors (packed into $X\in \mathbb{R}^{N\times D}$): 
$Q=XW_\mathrm{Q}$, $K=XW_\mathrm{K}$, $V=XW_\mathrm{V}$, using linear transformations $W_\mathrm{Q},W_\mathrm{K},W_\mathrm{V}$ with the constraint $k=N$, meaning that the attention is in between all the input vectors. \\
Finally, {Multi-head} self-attention layer (MSA) is defined by considering $h$ attention ``heads'', \textit{ie} $h$ self-attention functions  applied to the input. 
Each head provides a sequence of size $N\times d$. 
These $h$ sequences are rearranged into a $N \times dh$ sequence that is reprojected by a linear layer into $N \times D$. 

\paragraph{Transformer block for images.}
To get a full transformer block as in~\cite{Vaswani2017AttentionIA},  we add a Feed-Forward Network (FFN) on top of the MSA layer. This FFN is composed of two linear layers separated by a GeLu activation~\cite{Hendrycks2016GaussianEL}.
The first linear layer expands the dimension from $D$ to $4D$, and the second layer reduces the dimension from $4D$ back to $D$.
Both MSA and FFN are operating as residual operators thank to skip-connections, and with a layer normalization~\cite{ba2016layer}.

In order to get a transformer to process images, our work builds upon the ViT model~\cite{dosovitskiy2020image}.  It is a simple and elegant architecture that processes input images as if they were a sequence of input tokens. The fixed-size input RGB image is decomposed into a batch of $N$ patches of a fixed size of $16\times 16$ pixels ($N=14\times 14$). Each patch is projected with a linear layer that conserves its overall dimension $3 \times 16 \times 16=768$.\\
%
The transformer block described above is invariant to the order of the patch embeddings, and thus does not consider their relative position.
The positional information is incorporated as fixed~\cite{Vaswani2017AttentionIA} or trainable~\cite{ghering2017convseq2seq} positional embeddings. 
They are added before the first transformer block to the patch tokens, which are then fed to the stack of transformer blocks. 




\paragraph{The class token} is a trainable vector, appended to the patch tokens before the first layer, that goes through the transformer layers, and is then projected with a linear layer to predict the class.
This class token is inherited from NLP~\cite{devlin2018bert}, and departs from the typical pooling layers used in computer vision to predict the class. 
The transformer thus process batches of $(N + 1)$ tokens of dimension $D$, of which only the class vector is used to predict the output.
This architecture forces the self-attention to spread information between the patch tokens and the class token: at training time the supervision signal comes only from the class embedding, while the patch tokens are the model's only variable input.


\paragraph{Fixing the positional encoding across resolutions. }
\label{sec:position_encoding}

Touvron et al.~\cite{Touvron2019FixRes} show that it is  desirable to use a lower training resolution and fine-tune the network at the larger resolution. %
This speeds up the full training and improves the accuracy under prevailing data augmentation schemes. 
%
When increasing the resolution of an input image, we keep the patch size the same, 
%
% 
therefore the number $N$ of input patches does change.
Due to the architecture of transformer blocks and the class token, the model and classifier do not need to be modified to process more tokens.
In contrast, one needs to adapt the positional embeddings, because there are $N$ of them, one for each patch. 
Dosovitskiy et al.~\cite{dosovitskiy2020image}  interpolate the positional encoding when changing the resolution and demonstrate that this method works with the subsequent fine-tuning stage. 
